% =============================================================================
% PREÁMBULO LaTeX — Entrega P1, Bank Marketing
% -----------------------------------------------------------------------------
% Se inyecta en la plantilla numbats/report mediante la clave
% `include-in-header` del YAML de UTEC-Report1.qmd.
%
% Solo contiene ajustes tipográficos y de espaciado. Deliberadamente NO se
% redefine ningún comando de la plantilla ni se cargan paquetes que ésta ya
% incluya (amsmath, amssymb, booktabs, longtable, graphicx, hyperref), para
% evitar colisiones de opciones en la carga.
% =============================================================================

% --- Espaciado de tablas -----------------------------------------------------
% Las tablas de descriptivos son anchas y densas; un poco de aire vertical
% entre filas las hace legibles sin reducir más el tamaño de fuente.
\renewcommand{\arraystretch}{1.15}

% --- Control de viudas y huérfanas -------------------------------------------
% Evita que una línea suelta de un párrafo quede aislada al final o al inicio
% de página, algo frecuente en un documento con muchas figuras intercaladas.
\widowpenalty=10000
\clubpenalty=10000

% --- Figuras -----------------------------------------------------------------
% La plantilla permite que una figura ocupe una página entera; se relaja el
% umbral para que las figuras en rejilla (histogramas, boxplots) floten menos.
\renewcommand{\topfraction}{0.9}
\renewcommand{\bottomfraction}{0.8}
\renewcommand{\textfraction}{0.07}
\renewcommand{\floatpagefraction}{0.7}

% --- Notación estadística del informe ----------------------------------------
% Comandos propios para que la notación sea consistente en todo el documento
% y pueda cambiarse desde un único lugar.
\providecommand{\Cov}{\mathrm{Cov}}          % covarianza
\providecommand{\Var}{\mathrm{Var}}          % varianza
\providecommand{\RIC}{\mathrm{RIC}}          % rango intercuartílico
\providecommand{\corr}[2]{r_{#1#2}}          % coeficiente de correlación

% --- Resaltado de cifras clave -----------------------------------------------
% \cifra{11.27\%} destaca un valor numérico en el cuerpo del texto sin recurrir
% a negrita, que ya se usa para los términos conceptuales.
\providecommand{\cifra}[1]{\textbf{\texttt{#1}}}
